% ============================================================================
% Fichier  : partie5/P05_C25_exercices.tex
% Titre    : Les Exercices de Crise
% Auteur   : MINKA MI NGUIDJOÏ Thierry Emmanuel
% Version  : 1.0 -- novembre 2026
% ============================================================================
% Sources : ISO 22361:2022 ; ISO 22301:2019 (8.5) ; NIST SP 800-84 (2006).
% Exercice du 22 octobre 2026 de LiZbiZ : cas fictif.
% ============================================================================

\chapter{Les Exercices de Crise}
\label{chap:exercices-crise}
\chaptermark{Les exercices de crise}

\epigraph{%
    « On ne s'élève pas au niveau de ses attentes en crise~;\\
    on retombe au niveau de son entraînement. »%
}{Formule d'usage chez les instructeurs}

\niveauChapitre{Intermédiaire à expert}{Chapitres~\ref{chap:exercer},
\ref{chap:cellule} à~\ref{chap:communiquer}.}

\begin{boxobjectifs}
\begin{itemize}
    \item Concevoir un exercice de crise pour une cellule de direction.
    \item Écrire un scénario et un chronogramme d'injections réalistes.
    \item Animer, observer et préserver la sécurité de l'exercice.
    \item Conduire un débriefing à chaud et un débriefing à froid.
    \item Mesurer la progression d'une cellule d'un exercice à l'autre.
\end{itemize}
\end{boxobjectifs}

\bigskip

% ============================================================================
\section*{Question directrice : la cellule de crise de LiZbiZ ferait-elle mieux aujourd'hui qu'en mars~?}
% ============================================================================

Le 22~octobre 2026, Mme~MOUAFO réunit la direction de LiZbiZ pour un
premier exercice sur table de deux heures. Le scénario n'est pas
annoncé. La question n'est pas de « réussir », mais de mesurer, par
rapport au 6~mars, ce qui a progressé et ce qui manque encore.

\medskip
\textbf{Plan du chapitre.}
\begin{enumerate}
    \item Ce qu'entraîne un exercice de crise (section~\ref{sec:c25-objet}).
    \item Le scénario et les injections (section~\ref{sec:c25-scenario}).
    \item Animer et observer (section~\ref{sec:c25-animer}).
    \item Débriefer (section~\ref{sec:c25-debriefer}).
    \item Mesurer la progression (section~\ref{sec:c25-progression}).
\end{enumerate}

\separateur

% ============================================================================
\section{Ce qu'Entraîne un Exercice de Crise}
\label{sec:c25-objet}
% ============================================================================

Un exercice de continuité éprouve des procédures~; un exercice de crise
éprouve une \emph{organisation qui décide}~: se réunir vite, partager la
situation, décider sous incertitude, communiquer, tenir la main courante.
Il donne aux dirigeants l'expérience que les crises réelles, rares, ne
leur donnent pas (Chapitre~\ref{chap:decider}).

\begin{tableaucours}{L{4.6cm}L{7.6cm}}{%
    \tableauHeader{Capacité} &
    \tableauHeader{Objectif mesurable de l'exercice d'octobre}
}
Se réunir & Cellule réunie en moins de 30~minutes après l'alerte \\
Tenir la main courante & Main courante ouverte dans les 5~minutes \\
Communiquer & Premier message validé en moins de 60~minutes \\
Décider & Décisions conservatoires prises avant le premier point \\
Joindre & Tous les contacts utiles joints par l'annuaire de crise \\
\end{tableaucours}
\vspace{-0.8em}
\begin{center}{\small\itshape Tableau~25.1. Capacités entraînées et objectifs}\end{center}

% ============================================================================
\section{Le Scénario et les Injections}
\label{sec:c25-scenario}
% ============================================================================

Le scénario part du registre des risques. En octobre, Mme~MOUAFO choisit
un risque de la zone de contrôle de l'écosystème
(Chapitre~\ref{chap:ebios})~: une fuite de données chez un transporteur
sous-traitant.

\begin{tableaucours}{L{1.4cm}L{7.8cm}L{3.0cm}}{%
    \tableauHeader{Temps} &
    \tableauHeader{Injection} &
    \tableauHeader{Capacité visée}
}
T+0 & Un transporteur signale qu'un fichier de 3\,000~adresses de
destinataires de LiZbiZ est en vente sur un forum & Se réunir \\
T+15 & Un journaliste appelle pour confirmation & Communiquer \\
T+35 & Le transporteur ignore depuis quand le fichier a fuité & Décider
sous incertitude \\
T+50 & Un client majeur menace de suspendre son contrat & Parties
prenantes \\
T+70 & Le directeur général doit quitter la salle (suppléant) & Suppléance \\
T+90 & Le vendeur propose de « retirer » le fichier contre paiement &
Politique sur les rançons \\
\end{tableaucours}
\vspace{-0.8em}
\begin{center}{\small\itshape Tableau~25.2. Chronogramme de l'exercice du 22 octobre}\end{center}

\begin{boxalert}
\textbf{Violation typique.} Un scénario si dur que la cellule est
submergée et n'apprend rien, ou si doux qu'elle ne rencontre aucune
difficulté.

\smallskip
\textbf{Règle~:} doser~: chaque injection met une capacité à l'épreuve,
une seule à la fois au début, puis plusieurs. L'animateur peut accélérer
ou ralentir selon la réaction de la cellule.
\end{boxalert}

% ============================================================================
\section{Animer et Observer}
\label{sec:c25-animer}
% ============================================================================

\begin{tableaucours}{L{3.2cm}L{9.0cm}}{%
    \tableauHeader{Rôle} &
    \tableauHeader{Mission}
}
Directeur d'exercice & Garant du cadre et de la sécurité~; peut suspendre \\
Animateurs & Jouent les tiers (journaliste, client, transporteur), envoient
les injections \\
Observateurs & Notent, sur une grille, les heures et les faits liés à
chaque objectif \\
Joueurs & La cellule de crise, dans ses rôles réels \\
\end{tableaucours}
\vspace{-0.8em}
\begin{center}{\small\itshape Tableau~25.3. Les rôles d'un exercice de crise}\end{center}

Chaque message porte la mention « EXERCICE »~; un mot d'arrêt est convenu
pour le cas où un incident réel surviendrait pendant l'exercice. Aucun
appel n'est passé à un vrai journaliste ni à un vrai client.

% ============================================================================
\section{Débriefer}
\label{sec:c25-debriefer}
% ============================================================================

\begin{boxprocedure}
\textbf{Deux débriefings}
\begin{enumerate}
    \item \textbf{À chaud}, juste après l'exercice, trente minutes~: chaque
          joueur dit ce qui l'a aidé, ce qui l'a gêné, ce qu'il ferait
          autrement. On ne juge personne.
    \item \textbf{À froid}, une à deux semaines plus tard~: rapport fondé
          sur les grilles d'observation~; objectifs atteints ou non~;
          actions, responsables, échéances.
\end{enumerate}
\end{boxprocedure}

\begin{boxscenario}{Opération LIZBIZ. L'exercice du 22 octobre}
\textbf{Constatations} (grilles d'observation).
\begin{itemize}
    \item Cellule réunie en 25~minutes~; main courante ouverte à
          4~minutes.
    \item Premier message validé à 52~minutes.
    \item Décision de refuser le paiement conforme à la politique, prise en
          8~minutes.
    \item Suppléance~: Mme~MBALLA prend la direction, mais un numéro sur
          vingt-trois de l'annuaire de crise est faux.
    \item Personne n'est chargé d'appeler l'assureur~; l'appel n'est pas
          passé.
\end{itemize}

\medskip
\textbf{Analyse} (confiance élevée). Quatre objectifs sur cinq sont atteints
(contre aucun le 6~mars). Les deux écarts sont simples à corriger~: mise à
jour trimestrielle de l'annuaire, ajout de l'assureur à la fiche réflexe
(Chapitre~\ref{chap:cellule}).
\end{boxscenario}

% ============================================================================
\section{Mesurer la Progression}
\label{sec:c25-progression}
% ============================================================================

\begin{tableaucours}{L{5.2cm}R{3.4cm}R{3.4cm}}{%
    \tableauHeader{Indicateur} &
    \tableauHeader{6~mars (réel)} &
    \tableauHeader{22~octobre (exercice)}
}
Délai de réunion de la cellule & 3~h & 25~min \\
Délai d'ouverture de la main courante & Plus d'un jour & 4~min \\
Délai du premier message externe & 25~h~30 & 52~min \\
Contacts joints par l'annuaire & Pas d'annuaire & 22 sur 23 \\
\end{tableaucours}
\vspace{-0.8em}
\begin{center}{\small\itshape Tableau~25.4. Progression de la cellule de crise de LiZbiZ}\end{center}

\begin{boxalert}
\textbf{Violation typique.} Comparer un exercice préparé et une crise
réelle comme s'ils étaient équivalents.

\smallskip
\textbf{Règle~:} l'exercice se déroule un jour ouvré, à heure connue, avec
une cellule présente~; la crise de mars a éclaté un vendredi à l'aube, le
DG sur la route. La progression est réelle, mais elle doit être confirmée
par un exercice inopiné, hors des heures ouvrées.
\end{boxalert}

\begin{boxmetier}{Directeur de crise}
\textbf{Jouer son propre rôle.} M.~ETEME participe aux exercices dans son
rôle réel et accepte d'être mis en défaut~: un directeur général qui se
fait remplacer à l'exercice ne sera pas prêt le jour venu.
\end{boxmetier}

\separateur

\begin{boxresume}
\textbf{Ce qu'il faut retenir de ce chapitre~:}
\begin{itemize}
    \item \textbf{Un exercice de crise} éprouve une organisation qui
          décide, avec des objectifs mesurables.
    \item \textbf{Le scénario} vient du registre~; les injections sont
          dosées et datées.
    \item \textbf{Animateurs, observateurs, directeur d'exercice}~; mention
          « EXERCICE » et mot d'arrêt.
    \item \textbf{Débriefing à chaud} sans jugement, \textbf{à froid} fondé
          sur les observations.
    \item \textbf{La progression} se mesure d'un exercice à l'autre, et se
          confirme par un exercice inopiné.
\end{itemize}
\end{boxresume}

\begin{boxinfo}
\textbf{Points de contrôle de l'élève.} Je sais~: fixer les objectifs d'un
exercice de crise~; écrire un chronogramme d'injections~; répartir les
rôles~; conduire les deux débriefings~; mesurer une progression.
\end{boxinfo}

\bigskip

\begin{boxexo}[title={\ding{45}\quad Exercice~25.1. Les objectifs \niveaubase}]
Combien d'objectifs du Tableau~25.1 l'exercice d'octobre a-t-il atteints~?
Lesquels ne l'ont pas été~?
\end{boxexo}

\begin{boxexo}[title={\ding{45}\quad Exercice~25.2. Les rôles \niveaubase}]
Qui, dans l'exercice, joue le journaliste~? Qui peut suspendre l'exercice~?
\end{boxexo}

\begin{boxexo}[title={\ding{45}\quad Exercice~25.3. La progression \niveaubase}]
Calculez, pour chaque indicateur du Tableau~25.4, le gain entre mars et
octobre.
\end{boxexo}

\begin{boxexo}[title={\ding{45}\quad Exercice~25.4. Une injection \niveaumoyen}]
Ajoutez deux injections au chronogramme d'octobre pour éprouver la
communication interne.
\end{boxexo}

\begin{boxexo}[title={\ding{45}\quad Exercice~25.5. La grille \niveaumoyen}]
Construisez la grille d'observation de l'objectif « premier message validé
en moins de 60~minutes ».
\end{boxexo}

\begin{boxexo}[title={\ding{45}\quad Exercice~25.6. À chaud \niveaumoyen}]
Préparez les trois questions du débriefing à chaud et la règle qui
empêche d'en faire un procès.
\end{boxexo}

\begin{boxexo}[title={\ding{45}\quad Exercice~25.7. L'assureur \niveaumoyen}]
L'appel à l'assureur a été oublié. Corrigez la fiche réflexe de la
première heure (Chapitre~\ref{chap:cellule}).
\end{boxexo}

\begin{boxexo}[title={\ding{45}\quad Exercice~25.8. Inopiné \niveaumoyen}]
Concevez un exercice inopiné d'une heure, un samedi matin, et dites ce
qu'il mesure que l'exercice d'octobre ne mesurait pas.
\end{boxexo}

\begin{boxexo}[title={\ding{45}\quad Exercice~25.9. L'exercice de février \niveauexpert}]
Concevez l'exercice sur table « rançongiciel » de février 2027~: objectifs,
scénario, chronogramme de huit injections, rôles, sécurité.
\end{boxexo}

\begin{boxexo}[title={\ding{45}\quad Exercice~25.10. Le rapport à froid \niveauexpert}]
Rédigez le rapport à froid de l'exercice du 22~octobre.
\end{boxexo}

\bigskip

\begin{boxinfo}
\textbf{Pour aller plus loin.}
\begin{itemize}
    \item ISO~22361~\cite{iso22361_2022}~; NIST SP~800-84~\cite{nist80084}.
\end{itemize}
\end{boxinfo}

\begin{boxinfo}
\textbf{Transition.} Exercices et crises réelles ne servent que si l'on en
tire les leçons. Le Chapitre~\ref{chap:retex} traite du retour
d'expérience et de la mesure de la résilience.
\end{boxinfo}
